\chapter{Unit 5: Disaster Risk Reduction, Hazard Scales \& Sendai Framework}

\begin{formulabox}{Disaster Risk Formulas \& Hazard Measurement Scales}
\begin{itemize}[leftmargin=1.2em]
    \item \textbf{Universal Disaster Risk Formulation (UNDRR)}:
    \begin{equation}
    \text{Disaster Risk} = \frac{\text{Hazard} \times \text{Exposure} \times \text{Vulnerability}}{\text{Coping Capacity}}
    \end{equation}
    \item \textbf{Earthquake Magnitude \& Energy Release}:
    \begin{itemize}
        \item \textbf{Richter Local Magnitude ($M_L$)}:
        \begin{equation}
        M_L = \log_{10}(A) - \log_{10}(A_0)
        \end{equation}
        A unit increase in magnitude represents a $10$-fold increase in ground motion amplitude.
        \item \textbf{Gutenberg-Richter Energy Release Equation}:
        \begin{equation}
        \log_{10} E = 4.8 + 1.5 M \implies E = 10^{4.8} \times 10^{1.5 M}
        \end{equation}
        Each unit increase in Richter magnitude corresponds to a $10^{1.5} \approx 31.62$-fold increase in radiated seismic strain energy.
        \item \textbf{Modified Mercalli Intensity (MMI)}: Qualitative 12-point scale (I: Not felt to XII: Total destruction) evaluating ground shaking and structural damage.
    \end{itemize}
    \item \textbf{IMD Cyclone Classification (Maximum Sustained Wind Speeds - 3 min avg)}:
    \begin{itemize}
        \item Depression: $31\text{--}49\text{ km/h}$ ($17\text{--}27\text{ knots}$)
        \item Deep Depression: $50\text{--}61\text{ km/h}$ ($28\text{--}33\text{ knots}$)
        \item Cyclonic Storm: $62\text{--}88\text{ km/h}$ ($34\text{--}47\text{ knots}$)
        \item Severe Cyclonic Storm: $89\text{--}117\text{ km/h}$ ($48\text{--}63\text{ knots}$)
        \item Very Severe Cyclonic Storm: $118\text{--}166\text{ km/h}$ ($64\text{--}89\text{ knots}$)
        \item Extremely Severe Cyclonic Storm: $167\text{--}221\text{ km/h}$ ($90\text{--}119\text{ knots}$)
        \item Super Cyclonic Storm: $\ge 222\text{ km/h}$ ($\ge 120\text{ knots}$)
    \end{itemize}
\end{itemize}
\end{formulabox}

\begin{lawbox}{Disaster Management Act, 2005 \& Sendai Framework (2015--2030)}
\begin{itemize}[leftmargin=1.2em]
    \item \textbf{Three-Tier Institutional Architecture (DM Act, 2005)}:
    \begin{itemize}
        \item \textbf{National Level}: National Disaster Management Authority (NDMA) headed by the Prime Minister. National Executive Committee (NEC) headed by Home Secretary.
        \item \textbf{State Level}: State Disaster Management Authority (SDMA) headed by the Chief Minister.
        \item \textbf{District Level}: District Disaster Management Authority (DDMA) headed by the District Magistrate / Collector.
        \item \textbf{Specialist Agencies}: National Disaster Response Force (NDRF) - dedicated multi-disciplinary battalions; National Institute of Disaster Management (NIDM) for training and capacity building.
    \end{itemize}
    \item \textbf{The 4 Sendai Priorities for Action (2015--2030)}:
    \begin{enumerate}
        \item Priority 1: Understanding disaster risk.
        \item Priority 2: Strengthening disaster risk governance to manage disaster risk.
        \item Priority 3: Investing in disaster risk reduction for resilience.
        \item Priority 4: Enhancing disaster preparedness for effective response and to \textbf{``Build Back Better''} in recovery, rehabilitation and reconstruction.
    \end{enumerate}
    \item \textbf{The 7 Global Sendai Targets (A to G)}: Reduce mortality (A), reduce affected people (B), reduce direct economic loss (C), reduce infrastructure damage (D), increase national/local DRR strategies by 2020 (E), enhance international cooperation (F), increase access to multi-hazard early warning systems (G).
\end{itemize}
\end{lawbox}

\begin{memorytip}
Remember that each +1 increase on the Richter scale means $\approx 31.6$ times more energy release, and a +2 increase means $31.62^2 = 1000$ times more energy release!
\end{memorytip}
